\clearpage

\section{Numerical Sequences and Series}
\subsection*{Exercises}
\addcontentsline{toc}{subsection}{Exercises}

\begin{exercise}{11}
Suppose $a_n>0$, $s_n=a_1+\cdots+a_n$, and $\sum a_n$ diverges.
\begin{enumerate}
    \item[(a)] Prove that $\sum\frac{a_n}{1+a_n}$ diverges.
    \item[(b)] Prove that 
\[
\frac{a_{N+1}}{s_{N+1}} + \cdots + \frac{a_{N+k}}{s_{N+k}} \geq 1- \frac{s_{N}}{s_{N+k}}
\]
and deduce that $\sum \frac{a_n}{s_n}$ diverges.
    \item[(c)] Prove that
\[
\frac{a_n}{s_n^2} \leq \frac{1}{s_{n-1}}-\frac{1}{s_n}
\]
and deduce that $\sum \frac{a_n}{s_n^2}$ converges.
    \item[(d)] What can be said about
\[
\sum \frac{a_n}{1+na_n} \text{ and } \sum \frac{a_n}{1+n^2a_n} ?
\]
\end{enumerate}
\end{exercise}
\begin{proof}
\begin{enumerate}
    \item[(a)] Assume $a_n$ diverges. 
If $\{n_i\}$ is a sequence of positive integers with $n_1<n_2<n_3<\cdots$, then we have
\[
\sum_{i=1}^{n_k} \frac{a_i}{1+a_i} \geq \sum_{i=1}^k \frac{a_{n_i}}{1+a_{n_i}}.
\]

Since $a_n$ diverges, we have $0<s=\limsup_{n \to \infty} a_n \leq \infty$. 
Pick a subsequence $\{a_{n_i}\}_{i \in \N}$ such that $a_{n_i} \to s$ and let $k \to \infty$.
Then right side of the last inequality diverges, since $\frac{a_{n_i}}{1+a_{n_i}}$ does not converge to $0$.
Hence left hand side diverges.

Next, assume $a_n$ converges and let $a_n \to 0$. (Otherwise, $\sum \frac{a_n}{1+a_n}$ diverges.)
Then there exists an integer $N_1$ such that $a_n<1$ for $n \geq N_1$.

Suppose, for a contradiction, that $\sum \frac{a_n}{1+a_n}$ converges.
Let $\varepsilon>0$ be given. 
Choose an integer $N_2$ such that 
\[
m \geq n \geq N_2 \text{ implies } \sum_{k=n}^m\frac{a_k}{1+a_k}  \leq \varepsilon/2.
\]

Now, let $N= \max \{N_1,N_2\}$ and let $\alpha= \sup_{k \geq N} a_k$.
If $m \geq n \geq N$ then
\begin{align*}
    \frac{a_n}{1+a_n}+ \cdots +\frac{a_m}{1+a_m}  &\geq \frac{1}{1+\alpha}(a_n+ \cdots +a_m) \\
    &=\frac{1}{1+\alpha} \sum_{k=n}^m a_k,
\end{align*}
hence 
\[
\sum_{k=n}^m a_k \leq \frac{1+\alpha}{2} \varepsilon \leq \varepsilon.
\]
This contradicts our hypothesis that $\sum a_n$ diverges.

\item[(b)] Since $a_n>0$, we have
\begin{align*}
\frac{a_{N+1}}{s_{N+1}}+ \cdots +\frac{a_{N+k}}{s_{N+k}} &\geq \frac{a_{N+1}+ \cdots a_{N+k}}{s_{N+k}} \\
&=\frac{s_{N+k}-s_N}{s_{N+k}} \\
&= 1-\frac{s_N}{s_{N+k}}.
\end{align*}
for every positive integer $N$.
This yields the desired inequality.

Suppose $\sum \frac{a_n}{s_n}$ converges.
Choose an integer $N$ such that $m \geq n \geq N$ implies 
\[
\sum_{k=n}^m \frac{a_k}{s_k} \leq 1/2.
\]
Hence we have
\[
1-\frac{s_N}{s_{N+k}} \leq \frac{a_{N+1}}{s_{N+1}} + \cdots + \frac{a_{N+k}}{s_{N+k}}
\]
for fixed $N$ and for all $k \in \N$.
Letting $k \to \infty$, we obtain $1 \leq \frac{1}{2}$, a contradiction.

\item[(c)] Observe that
\begin{align*}
    \frac{a_n}{s_n^2} &\leq \frac{a_n}{s_{n-1}s_n} \\
    &= \frac{a_n}{s_n-s_{n-1}}\left(\frac{1}{s_{n-1}}-\frac{1}{s_n}\right) \\
    &=\frac{1}{s_{n-1}}-\frac{1}{s_n}.
\end{align*}
Taking summation from $n=2$ of both sides and letting $n \to \infty$, by telescoping, we conclude that $\sum \frac{a_n}{s_n^2}$ converges.

\item[(d)] We first prove that the latter one converges.
Observe that
\[
\frac{a_n}{1+n^2a_n} \leq \frac{a_n}{n^2a_n}=\frac{1}{n^2},
\]
which implies that $\sum \frac{a_n}{1+n^2a_n}$ converges.

However, there is nothing that we can say about the former one.
For instance, if we put $a_n=\frac{1}{n}$ then $\sum \frac{a_n}{1+na_n}=\sum \frac{1}{2n}$ diverges. 
But if we take $a_{m^2}=1/m$ for $m=1,2,3, \ldots$, while $a_n=2^{-n}$ if $n$ is not a square, then the sum converges.\footnote{The last example is from this: https://math.stackexchange.com/questions/237691/baby-rudin-chapter-3-exercise-11d}
\end{enumerate}
\end{proof}

\begin{exercise}{20}
suppose $\{p_n\}$ is a Cauchy sequence in a metric space $X$, and some subsequence $\{p_{n_i}\}$ converges to a point $p \in X$.
Prove that the full sequence $\{p_n\}$ converges to $p$.
\end{exercise}
\begin{proof}
Fix $\varepsilon>0$.
Since $\{p_n\}$ is Cauchy, there exists an integer $N$ such that $d(p_n,p_m)<\varepsilon/2$ if $n \geq N$ and $m \geq N$.
Since a subsequence $\{p_{n_i}\}$ converges to $p$, choose an integer $k \geq N$ such that $d(p_{n_k},p)<\varepsilon/2$.
Hence, $n \geq N$ implies
\begin{align*}
d(p_n,p) &\leq d(p_n, p_{n_k})+d(p_{n_k}, p) \\
&<\varepsilon/2+\varepsilon/2=\varepsilon.
\end{align*}
This is the desired result.
\end{proof}